\documentclass[11pt]{article}

\usepackage{hw_style}

% Set due date for this homework (updated to 2025)
\setduedate{11}{14}

% Setup homework number and header
\hwnumber{6}

\begin{document}

\begin{hwproblem}{1}
证明对于矩形波,
$f(x)=\left\{ \begin{array}{ll}-1, & x<0 \\ +1, &  x>0\end{array}\right.$, 其傅里叶级数为
$$
f(x)=\frac{4}{\pi} \left[  \sin \pi x+\frac{1}{3} \sin 3 \pi+\frac{1}{5 } \sin 5 x+\cdots \right].
$$
\end{hwproblem}


\begin{hwproblem}{2}
计算矩形波$$
\begin{aligned}
& f(x)=0, \quad-\pi<x<0, \\
& f(x)=h, \quad 0<x<\pi .
\end{aligned}
$$
的傅里叶级数展开.注意观察级数展开在间断点的取值是否为预期.
\end{hwproblem}

\vspace{0.5cm}
  
\begin{hwproblem}{3}
    证明傅里叶变换$\mathcal{F}[f(x)] = F(\omega)$ 满足以下定理.
    \begin{enumerate}
        % \item 导数定理
        % $$
        %     \mathcal{F} [f'(x)] = \ii \omega F(\omega)
        % $$
      
        \item 积分定理
        $$
            \mathcal{F} [ \int^{x} f(x) dx ] = \frac{1}{\ii \omega} F(\omega),
        $$
        \item 位移定理
        $$
            \mathcal{F} [ e^{\ii \omega_0 x} f(x) ] = F(\omega - \omega_0),
        $$
        \item 卷积定理(定义见书本)
        $$
            \mathcal{F} [f_1(x)\star f_2(x) ] = F_1(\omega) F_2(\omega).
        $$
      \end{enumerate}
\end{hwproblem}

\vspace{0.5cm}

\begin{hwproblem}{4}
    对于屏蔽效应的描述可以使用Yukawa势来表示
$$
V(\mathbf{r}) = \frac{e^{-\alpha r}}{r},
$$
其中$\alpha$描述屏蔽的强弱.
\begin{enumerate}
  \item 试在球坐标系下求出三维傅里叶变换;
  \item 当$\alpha=Z$时, 函数$e^{-Zr}$可描述类氢原子的$1s$轨道, 对其傅里叶变换关于$Z$求偏导
  $-\frac{\partial}{\partial Z} \mathcal{F} \left[ \frac{e^{-Z r}}{r}\right]$;
  \item 若是三维高斯形式$V(\mathbf{r}) = e^{-\alpha r^2}$,求其傅里叶变换.
\end{enumerate}
\end{hwproblem}

\vspace{0.5cm}
  
  
\begin{hwproblem}{5}
    \begin{enumerate}
        \item 验证$f(x)=x^{-1 / 2}$ 在Fourier傅里叶余弦和正弦变换下互易;即
        $$
        \begin{aligned}
        & \sqrt{\frac{2}{\pi}} \int_0^{\infty} x^{-1 / 2} \cos {x t} \, dx=t^{-1 / 2} \\
        & \sqrt{\frac{2}{\pi}} \int_0^{\infty} x^{-1 / 2} \sin {x t} \, dx=t^{-1 / 2}
        \end{aligned}
        $$
        \item 使用上述结果计算菲涅耳积分(Fresnel integrals)
        $$
        \int_0^{\infty} \cos \left(y^2\right) d y \text { , } \int_0^{\infty} \sin \left(y^2\right) d y
        $$
      \end{enumerate}
\end{hwproblem}

\vspace{0.5cm}

\hwrequirements

\end{document}